\section{Related work}\label{sec:related}

\subsection{Household load forecasting with deep learning}
Household load is volatile and dominated by occupant behaviour. Kong et al.~\cite{kong2019stlf} showed that LSTMs
outperform shallow and statistical baselines for individual residential customers. Shi et al.~\cite{shi2018pdrnn}
showed that pooling the data of neighbouring households reduces over-fitting, the effect that motivates
collaborative training in the first place. Later work added embeddings, bidirectional recurrences and multi-task
sharing~\cite{aurangzeb2024,xiao2023}, as well as clustering of consumption profiles~\cite{dab2023}. These gains assume
the data are pooled.

\subsection{Federated household load forecasting}
FL for residential forecasting has grown quickly~\cite{fernndez2022,fekri2022,gholizadeh2022,yang2023,liu2023,fekri2023,
qu2023,wang2023,khalil2023,widmer2023,dogra2023,sarmento2024,abdulla2024,singh2024,manzoor2024,mao2024,huang2024,
mostafanezhad2025,kaur2025,shao2026,das2026,wang2026,tibermacine2026,he2026,lu2024fds,husnoo2024,he2023hfl,
duttagupta2023,rahman2025multihop}.
On the question that matters most here, whether federation pays off at all, these studies disagree. The closest
to ours is that of Briggs et al.~\cite{briggs2022}, who trained a two-layer LSTM with FedAvg on 100 LCL households.
Plain FL was 7.6\,\% worse in mean RMSE than localised learning; clustering reduced the shortfall to 4.7\,\%, and
local fine-tuning made the federated models better than the local ones. Gholizadeh and Musilek~\cite{gholizadeh2022}
also found local learning best on average for individual Canadian homes, while Fekri et al.~\cite{fekri2022} found
FedAvg comparable to, or better than, per-meter models on 19 Ontario homes. Personalised and clustered
variants~\cite{qu2023,widmer2023,dogra2023,das2026,wang2026,lu2024fds} are almost always motivated by heterogeneity
among households.
None of these studies varies heterogeneity while holding the households and the test set fixed, so the share of the
federated gap caused by heterogeneity, rather than by data fragmentation or the training protocol, has not been
measured. We revisit the LCL comparison of Briggs et al. under a controlled protocol, with per-household paired
statistics, and add that decomposition.

\subsection{Communication cost}
Communication efficiency was a founding motivation of FL~\cite{mcmahan2017fedavg,konecny2016comm}.
In load forecasting the evidence on it points both ways. Taïk and Cherkaoui~\cite{taik2020} estimated analytically
that FL saves 90--97\,\% of the network load of uploading data, and Savi and Olivadese~\cite{savi2021} reported up to
50 times less exchanged data at sub-hourly resolution. Briggs et al.~\cite{briggs2020ccai}, on the other hand,
argued without measurement that exchanging weight vectors costs more bandwidth than sending half-hourly readings over
constrained smart-meter links.
Most FL forecasting studies count rounds rather than bytes~\cite{fekri2022}, or report payload per round without
comparing it with the data~\cite{bose2023pl}. We measure upload and download bytes per round, the bytes needed to
reach a given accuracy, and the raw-data volume they replace.

\subsection{Non-IID data and proximal regularisation}
Statistical heterogeneity slows and biases FedAvg~\cite{zhao2018noniid,hsieh2020quagmire,kairouz2021advances}.
The remedies proposed include proximal regularisation (FedProx~\cite{li2020fedprox}), control variates
(SCAFFOLD~\cite{karimireddy2020scaffold}), dynamic regularisation (FedDyn~\cite{acar2021feddyn}), normalised
averaging (FedNova~\cite{wang2020fednova}), adaptive server optimisers~\cite{reddi2021fedopt},
personalisation~\cite{fallah2020perfedavg,li2021ditto} and fairness-aware objectives~\cite{li2020qffl}.
Heterogeneity is usually induced with a Dirichlet split over class labels~\cite{hsu2019dirichlet}. Regression on
smart-meter data has no class labels, so we use a Dirichlet split over load-profile clusters and also measure the
resulting heterogeneity directly.

Choosing the proximal coefficient remains ad hoc. The FedProx paper adapts $\mu$ from the training
loss~\cite{li2020fedprox}. HAPI-FedProx adapts the proximal term from a local-versus-global performance
index~\cite{mhou2025hapi}. Yu et al.~\cite{yu2025pflminimax} characterise the optimal regularisation strength as a
function of heterogeneity in theory. FedAdp~\cite{wu2021fedadp} uses the geometry of client updates to set
aggregation weights. The update-disagreement ratio used by HA-FedProx is a bounded, update-space estimate of
FedProx's $B$-dissimilarity, and it belongs to the same family as gradient diversity~\cite{yin2018diversity}. To our
knowledge it has not been used to schedule $\mu$, and its behaviour on real energy data has not been examined.

\subsection{Privacy}
FL keeps raw data local, but updates can leak training data through gradient inversion~\cite{zhu2019leakage,
geiping2020inverting}. We therefore claim data locality only, not formal privacy.
